\BeleskeID{08-s061}
% element: 08-s061:e04
% element: 08-s061:notes
Prvi interval određuje napon koji kondenzator dostigne neposredno pre $t_2$. U trenutku druge promene napon kondenzatora je kontinualan, pa ta vrednost postaje početni izlaz drugog intervala. Nova konstantna pobuda iznosi $V_{OL}$ i zato je to i novo krajnje stanje.

Neposredna zamena početne i krajnje vrednosti daje međukorak
\[u_o(t)=V_{OL}+\left(V_{OH}+(V_{OL}-V_{OH})e^{-\frac{t_2-t_1}{\tau}}-V_{OL}\right)e^{-\frac{t-t_2}{\tau}},\quad t\geq t_2.\]
Izdvajanjem $V_{OH}-V_{OL}$ dobija se završni izraz. Rezultat je isti kao pri superpoziciji; postupak po intervalima jasnije pokazuje prenošenje stanja kroz trenutak druge promene.
